% Retrievalprestatie instellingsbesluit-matcher (DV1).
% Bron: instellingsbesluit_data/current_no_vlam_top25_20260512_validation_metrics.json
%        + MRR berekend uit current_no_vlam_top25_20260512_gold_link_ranks.csv (0,683 over 151 gold links).
% Run: no-VLAM, deterministische tekstuele retrieval + heuristische herrangschikking, 12 mei 2026.

\subsubsection{Retrievalprestatie van de instellingsbesluit-matcher.}
De instellingsbesluit-matcher werkt in twee fasen: een retrievallaag zoekt
kandidaat-documenten die mogelijk het instellingsbesluit van een
adviescollege(-fase) vormen, waarna de VLAM-poort die kandidaten inhoudelijk
beoordeelt (Tabel~\ref{tab:dv1-matchers}). De hier gerapporteerde evaluatie
betreft uitsluitend de retrievallaag: de vraag of het juiste instellingsbesluit
\"uberhaupt in de kandidaatlijst verschijnt, nog v\'o\'or de VLAM-beoordeling.
Voor deze meting stond VLAM uit en draaide alleen de deterministische, tekstuele
retrieval met heuristische herrangschikking, zodat de gerapporteerde recall de
prestatie van de zoeklaag op zichzelf weergeeft.

\noindent Als grondwaarheid dienden eerder handmatig gevalideerde canonieke
koppelingen tussen een adviescollege-fase en het bijbehorende instellingsbesluit
(Staatsblad of Staatscourant). Het ging om 151 canonieke koppelingen over 145
adviescollege-fasen. Net als bij de kabinetsreactie-matcher gaat het dus om
recall t\'egen bekende, gevalideerde paren, niet om volledige recall over alle
werkelijke koppelingen (zie Tabel~\ref{tab:dv1-beperkingen}).

\noindent Het juiste instellingsbesluit staat in 63\,\% van de gevallen direct op
de eerste positie, in 76\,\% binnen de top vijf en in 79\,\% in de volledige
kandidaatlijst; de mediane rangpositie van een gevonden document is 1 (maximaal
44). De mean reciprocal rank (MRR), het gemiddelde van de reciproque
rangposities, bedraagt 0,68, wat aansluit bij een juist document dat doorgaans
bovenaan de lijst staat (Tabel~\ref{tab:dv1-instelling-retrieval}). In 117 van de
151 gevallen (78\,\%) belandde het juiste document in de shortlist die normaal
aan VLAM wordt voorgelegd. De precisie van de retrievallaag is bewust laag --
slechts 5,7\,\% van de shortlist-kandidaten komt overeen met een gevalideerde
koppeling -- passend bij een brede zoekstap die op recall is ingericht en de
uiteindelijke selectie aan de VLAM-poort overlaat.

\noindent De prestatie verschilt sterk per type college: voor eenmalige colleges
zit het juiste besluit in 88\,\% van de gevallen in de shortlist, voor tijdelijke
colleges in 73\,\% en voor permanente colleges in slechts 55\,\%. Permanente
colleges zijn vaker via wetten, oudere Staatsblad-documenten of bredere
wettelijke kaders ingesteld, wat de retrieval bemoeilijkt. De schattingen zijn
omgeven met onzekerheid, en de grondwaarheid bestaat uit eerder gevalideerde
koppelingen, zodat lacunes daarin doorwerken in de cijfers.

\begin{table}[htbp]
\centering
\caption{Retrievalprestatie van de instellingsbesluit-matcher tegen handmatig
gevalideerde canonieke koppelingen ($N = 151$ koppelingen over 145
adviescollege-fasen), met 95\%-betrouwbaarheidsintervallen volgens Wilson. Alleen
de deterministische tekstuele retrieval stond aan; VLAM was uitgeschakeld, zodat
de cijfers de zoeklaag op zichzelf weergeven.}
\label{tab:dv1-instelling-retrieval}
\footnotesize
\begin{tabular}{@{}lcc@{}}
\toprule
Metriek & Waarde & 95\,\%-BI \\
\midrule
Recall@1               & 0,63 & 0,55--0,70 \\
Recall@5               & 0,76 & 0,68--0,82 \\
Recall@10              & 0,78 & 0,70--0,83 \\
Recall@25              & 0,78 & 0,71--0,84 \\
Candidate recall       & 0,79 & 0,72--0,85 \\
VLAM-shortlist recall  & 0,78 & 0,70--0,83 \\
MRR                    & 0,68 & --- \\
\bottomrule
\end{tabular}
\end{table}
